%=============================================================================
% Work Sheets (Appendix)
%=============================================================================
\chapter{Work Sheets}
\label{ch:worksheets}

%=============================================================================
\section{Work Sheet 5 --- Functions}
\label{ws:functions}
%=============================================================================

\begin{importantbox}{Instructions}
For each question below, write a complete C++ program.  Where a
\emph{function} is requested, implement the logic inside a
user-defined function and call it from \texttt{main()}.
\end{importantbox}

\begin{enumerate}[leftmargin=*, label=\textbf{Q\arabic*.}]

\item Write a program that reads 20 characters and uses a function to
      count and return how many of them are \textbf{uppercase} letters.

\item Write a function that receives a distance expressed in
      \textbf{feet and inches} and converts it to \textbf{meters}.\\
      \textit{Hint:} 1 foot = 12 inches, 1 inch = 2.54\,cm.

\item Write a function that receives a total number of \textbf{seconds}
      and converts it to the format \texttt{hr:mn:sec}.  For example,
      3661 seconds $\rightarrow$ \texttt{1:1:1}.

\item Write a function that takes an integer and determines whether it
      is \textbf{odd} or \textbf{even}.

\item Write a function that computes the \textbf{permutation}
      $P(n) = n!$ of a given positive integer~$n$.

\item Write a function that receives a student's mark and returns the
      corresponding \textbf{grade point} according to the following
      scale:

      \begin{center}
      \renewcommand{\arraystretch}{1.2}
      \begin{tabular}{@{} c c @{}}
      \toprule
      \textbf{Mark Range} & \textbf{Grade Point} \\
      \midrule
      90--100 & 4 \\
      80--89  & 3 \\
      70--79  & 2 \\
      60--69  & 1 \\
      Below 60 & 0 \\
      \bottomrule
      \end{tabular}
      \end{center}

\item Write a program that uses a function to compute and print the
      first $n$ terms of the \textbf{Fibonacci} sequence:\\
      $0,\ 1,\ 1,\ 2,\ 3,\ 5,\ 8,\ 13,\ \ldots$

\item Write a function that receives an integer $n$ and returns its
      \textbf{factorial} ($n!$).

\item Using a factorial function, evaluate
      \[
          z = \frac{x! - y!}{(x - y)!}
      \]
      for user-supplied values of $x$ and $y$ (where $x \geq y$).

\item Write a function that receives a \textbf{year} and determines
      whether it is a \textbf{leap year}.\\
      \textit{Rule:} A year is a leap year if it is divisible by~4,
      except for century years which must also be divisible by~400.

\item Write a program that reads two integers $x$ and $y$ and uses the
      library function \texttt{pow()} to compute $x^{y}$.

\item Write a function that receives an integer and returns its
      \textbf{reverse}.  For example, $765432 \rightarrow 234567$.

\item Write a program that uses functions to read the marks of
      \textbf{8~students}, then compute and display their \textbf{sum}
      and \textbf{average}.

\item Write a function that receives a character and \textbf{converts
      its case} --- uppercase to lowercase or vice versa --- and returns
      the result.

\item Write a \textbf{recursive} function that computes $n^{p}$
      (the power of $n$ raised to $p$) without using
      \texttt{pow()}.\\
      \textit{Hint:} $n^{p} = n \times n^{p-1}$, with base case
      $n^{0} = 1$.

\end{enumerate}

\newpage
%=============================================================================
\section{Work Sheet 6 --- Arrays}
\label{ws:arrays}
%=============================================================================

\begin{importantbox}{Instructions}
For each question below, write a complete C++ program.  Use arrays as
specified; you may use functions where appropriate.
\end{importantbox}

\begin{enumerate}[leftmargin=*, label=\textbf{Q\arabic*.}]

\item Write a program that reads an array of $n$ integers and checks
      whether the array is \textbf{sorted in ascending order}.

\item Write a program that reads an array of $n$ integers and counts
      the number of elements that are equal to \textbf{zero}.

\item Write a program that reads two integer arrays of the same size
      and produces a third array that is their element-wise
      \textbf{sum}.

\item Write a program that reads an integer array and creates a new
      array in which every element is \textbf{multiplied by~2}.

\item Write a program that reads the daily temperatures for one week
      (7 values) into an array, then computes and prints the
      \textbf{average temperature}.

\item Write a program that reads two sorted arrays and \textbf{merges}
      them into a single sorted array.

\item Write a program that reads a $3 \times 4$ matrix and computes
      the \textbf{sum of each column}.

\item Write a program that reads a square matrix of order $n$ and
      computes the \textbf{sum of the main diagonal} elements.

\item Write a program that reads a square matrix of order $n$ and
      computes the \textbf{sum of the secondary diagonal} elements.

\item Write a program that reads a $4 \times 5$ matrix and
      \textbf{searches} for a user-supplied value, reporting its
      position (row and column) if found.

\item Write a program that reads a one-dimensional array of 12 elements
      and stores them into a $3 \times 4$ two-dimensional array
      (\textbf{1-D to 2-D conversion}).

\item Write a program that reads a character array and \textbf{splits}
      it into two arrays: one containing the letters and the other
      containing the digits.

\item Write a program that reads an integer array and replaces every
      element that is both \textbf{even} and \textbf{positive} with
      the value~\textbf{10}.  Print the array before and after the
      replacement.

\item Write a program that reads a square matrix and replaces every
      \textbf{even} element on the \textbf{main diagonal} with~0.
      Print the matrix before and after.

\item Write a program that reads a $4 \times 5$ matrix and finds the
      \textbf{minimum element} in the entire matrix, printing its value
      and position (row and column).

\item Write a program that reads a $3 \times 3$ matrix and
      \textbf{exchanges} (swaps) two user-specified \textbf{rows}.
      Print the matrix before and after the swap.

\item Write a program that reads a $3 \times 3$ matrix and
      \textbf{exchanges} (swaps) a user-specified \textbf{row} with a
      user-specified \textbf{column}.  Print the matrix before and
      after.

\item Write a program that reads an array of $n$ integers and performs
      a \textbf{left rotation} by one position.  For example,
      $\{1,2,3,4,5\} \rightarrow \{2,3,4,5,1\}$.

\item Write a program that reads an array of $n$ integers and
      \textbf{deletes} the element at a user-specified position,
      shifting subsequent elements to fill the gap.

\item Write a program that reads an array of $n$ integers and sorts it
      in \textbf{ascending} order.

\item Write a program that reads an array of $n$ integers and sorts it
      in \textbf{descending} order.

\item Write a program that reads an array of $n$ integers and computes
      the \textbf{average of the even} numbers only.

\item Write a program that reads a $3 \times 4$ matrix~$A$ and creates
      a one-dimensional array~$B$ of size~3 such that each element
      $B[i]$ is the \textbf{sum of row~$i$} of~$A$.

      \textit{Example:}
\begin{verbatim}
A = | 1  2  3  4 |    B = | 10 |
    | 5  6  7  8 |        | 26 |
    | 9 10 11 12 |        | 42 |
\end{verbatim}

\item Write a program that reads a square matrix of order $n$ and
      finds the \textbf{greatest element on the secondary diagonal}.

\item Write a program that creates and prints each of the following
      arrays (the user supplies the size~$n$):

      \begin{enumerate}[label=(\alph*)]
        \item An array containing the first $n$ \textbf{even} numbers:
              $\{2, 4, 6, 8, \ldots\}$
        \item An array containing the first $n$ \textbf{odd} numbers:
              $\{1, 3, 5, 7, \ldots\}$
        \item An array whose elements follow the pattern
              $\{1, 2, 4, 8, 16, \ldots\}$ (powers of~2).
      \end{enumerate}

\end{enumerate}

\newpage
%=============================================================================
\section{Work Sheet 7 --- Strings}
\label{ws:strings}
%=============================================================================

\begin{importantbox}{Instructions}
For each question below, write a complete C++ program using C-style
strings (\texttt{char} arrays).  Make use of the string I/O and
library functions discussed in Chapter~\ref{ch:strings}.
\end{importantbox}

\begin{enumerate}[leftmargin=*, label=\textbf{Q\arabic*.}]

\item Write a program that reads a string from the user, prints the
      complete string, and then prints its characters in
      \textbf{reverse order} (one character per line).

      \textit{Example:}
\begin{verbatim}
Input:  Hello
Output: Hello
        o
        l
        l
        e
        H
\end{verbatim}

\item Write a program that reads a string and \textbf{swaps the case}
      of every letter --- uppercase letters become lowercase and vice
      versa.  Non-letter characters remain unchanged.

      \textit{Example:}
\begin{verbatim}
Input:  Hello World 123
Output: hELLO wORLD 123
\end{verbatim}

\item Write a program that reads a full sentence (using
      \texttt{cin.getline()}) and prints each \textbf{word on a
      separate line}.  Assume words are separated by single spaces.

      \textit{Example:}
\begin{verbatim}
Input:  Structured Programming in C++
Output: Structured
        Programming
        in
        C++
\end{verbatim}

\item Write a program that demonstrates the use of the following
      \textbf{I/O functions}:
      \begin{itemize}
        \item \texttt{cin.getline()} --- to read a full line,
        \item \texttt{cin.get()} --- to read a single character,
        \item \texttt{cin.ignore()} --- to discard remaining input,
        \item \texttt{cin.putback()} --- to push a character back,
        \item \texttt{cout.put()} --- to output a single character.
      \end{itemize}
      Display appropriate messages so the output clearly shows which
      function is being used.

\item Write a program that demonstrates the use of the following
      \textbf{string library functions}:
      \begin{itemize}
        \item \texttt{strlen()} --- to find the length of a string,
        \item \texttt{strcpy()} --- to copy one string into another,
        \item \texttt{strcat()} --- to concatenate two strings,
        \item \texttt{strcmp()} --- to compare two strings.
      \end{itemize}
      Read two strings from the user and apply each function, printing
      the result after every operation.

\end{enumerate}
